\chapter*{คำนำ}
\addcontentsline{toc}{chapter}{คำนำ}

จุลชีววิทยาสิ่งแวดล้อม (Environmental Microbiology) เป็นศาสตร์แขนงสำคัญที่ศึกษาความหลากหลาย บทบาทหน้าที่ กิจกรรมทางชีวเคมี และปฏิสัมพันธ์ของกลุ่มจุลินทรีย์ในระบบนิเวศธรรมชาติ ไม่ว่าจะเป็นในดิน น้ำ อากาศ ตลอดจนสภาพแวดล้อมที่มีความสุดขั้ว โดยเฉพาะอย่างยิ่งในมิติด้านทรัพยากรน้ำ ซึ่งเป็นปัจจัยพื้นฐานที่ค้ำจุนทั้งสิ่งมีชีวิต กิจกรรมทางการเกษตร และการพัฒนาอุตสาหกรรม การทำความเข้าใจโครงสร้างประชากรจุลินทรีย์ การตรวจติดตามตัวบ่งชี้ทางชีวภาพ และกลไกการย่อยสลายทางชีวภาพ จึงเป็นกุญแจสำคัญในการจัดการและฟื้นฟูคุณภาพสิ่งแวดล้อมอย่างยั่งยืน

ตำราวิชาการเรื่อง \textbf{``จุลชีววิทยาสิ่งแวดล้อมและการวิเคราะห์คุณภาพน้ำ (Environmental Microbiology and Water Quality Analysis)''} เล่มนี้ ได้รับการพัฒนาขึ้นโดยหลอมรวมองค์ความรู้ด้านจุลชีววิทยาบริสุทธิ์และจุลชีววิทยาประยุกต์เข้ากับระเบียบวิธีวิจัยและคู่มือปฏิบัติการเชิงลึก เนื้อหาภายในเล่มแบ่งออกเป็น 12 บทเรียนหลัก ครอบคลุมตั้งแต่หลักความปลอดภัยในห้องปฏิบัติการทางชีวภาพ การสำรวจการแพร่กระจายของจุลินทรีย์ในสิ่งแวดล้อม เทคนิคการย้อมสีและสัณฐานวิทยาของเซลล์ การแยกเชื้อบริสุทธิ์ อาหารเลี้ยงเชื้อ จลนศาสตร์การเจริญ จุลินทรีย์ไร้อากาศ เมแทบอลิซึมและการทดสอบทางชีวเคมี การควบคุมจุลินทรีย์ ตลอดจนการประยุกต์ใช้นวัตกรรมไบโอบล็อก (Bioblock) และจุลินทรีย์ตรึงรูปในการบำบัดน้ำเสียชุมชนและอุตสาหกรรม

ผู้เรียบเรียงหวังเป็นอย่างยิ่งว่า ตำราเล่มนี้จะเป็นประโยชน์ต่อนักศึกษาในสาขาวิชาจุลชีววิทยา วิทยาศาสตร์สิ่งแวดล้อม ดิจิทัลการเกษตรและสิ่งแวดล้อม คณาจารย์ นักวิจัย ตลอดจนนักวิชาการสิ่งแวดล้อม เพื่อนำไปใช้เป็นคู่มือประกอบการเรียนการสอน การวิจัย และการประยุกต์ใช้เทคโนโลยีชีวภาพเพื่อการอนุรักษ์สิ่งแวดล้อมและคุณภาพน้ำของประเทศสืบไป

\vspace{1.5cm}
\begin{flushright}
    \textbf{ผู้ช่วยศาสตราจารย์ ดร.จิรภัทร จันทมาลี}\\[2pt]
    \textbf{ผู้ช่วยศาสตราจารย์ ดร.ชีวะ ทัศนา}\\[4pt]
    คณะวิทยาศาสตร์และเทคโนโลยี มหาวิทยาลัยราชภัฏรำไพพรรณี\\[2pt]
    พุทธศักราช 2569
\end{flushright}
\clearpage
